\chapter{Paracetamol}

\medskip
\noindent\textit{\textbf{Under review.} This chapter is an AI-assisted educational draft; it is not a substitute for professional medical advice.}

\section*{Definition and Overview}
paracetamol is a non-opioid analgesic and antipyretic. Excess dosing can cause serious, initially silent liver injury.

\section*{Clinical Features}
pain relief or fever reduction at appropriate doses; overdose may cause nausea, sweating, abdominal discomfort, or no early symptoms.

\section*{When to Seek Medical Care}
Seek medical review for new, persistent, worsening, or function-limiting symptoms. Call 000 for severe breathing difficulty, collapse, sudden neurological change, severe chest or abdominal pain, or any rapidly worsening illness.

\section*{Diagnosis and Investigations}
establish dose and timing, check all combination products, and obtain urgent poison or emergency advice after an overdose; blood testing guides antidote treatment.

\section*{Management}
use only the labelled or prescribed dose, avoid duplicate products, and seek immediate advice after exceeding the dose. N-acetylcysteine is used in hospital when indicated.

\section*{Complications}
Complications depend on the underlying cause and may include progression, effects on other organs, treatment adverse effects, or reduced quality of life.

\section*{Prognosis and Self-Care}
Outcome depends on the cause, severity, complications, and response to treatment. Follow the care plan, keep appointments, avoid known triggers, and seek help if symptoms worsen.

\clearpage
